\chapter*{Exposé XXIII\\Reductive Groups: Uniqueness of Pinned Groups}
\addcontentsline{toc}{chapter}{Exposé XXIII. Uniqueness of Pinned Groups}
\begin{center}
\textit{by M. Demazure}
\end{center}

The aim of this Exposé is to prove the uniqueness theorem
(Theorem~4.1). It was proved by Chevalley in the case of an
algebraically closed field; the reduction-to-rank-two method used here
is likewise due to Chevalley (see \textit{Bible}, Exposés~23 and~24).
Along the way, we obtain an explicit description of reductive groups
by generators and relations (3.5).

\SGANextSectionTarget{sga3:xxiii:section:1}\SGASection{1}{Pinnings}

\SGARefTarget{sga3:XXIII:1.1}
\SGAStatement{Definition 1.1}{Let \(S\) be a scheme and let
\((G,T,M,R)\) be an \(S\)-split group (XXII, \SGARefLink{sga3:xxii:def:1.13}{1.13}). A \emph{pinning}
of this split group is the datum of a system \(\Delta\) of simple
roots of \(R\), together with, for each \(\alpha\in\Delta\), a section
\[
  X_\alpha\in\Gamma(S,\mathfrak g^\alpha)^\times .
\]}

In other words, a pinning of the reductive group \(G\) over the
nonempty scheme \(S\) is the datum of:
\begin{enumerate}
  \item[(i)] a maximal torus \(T\);
  \item[(ii)] an abelian group \(M\) and an isomorphism
    \(T\simeq D_S(M)\);
  \item[(iii)] a system of roots \(R\) of \(G\) with respect to \(T\);
  \item[(iv)] a system \(\Delta\) of simple roots of \(R\);
  \item[\SGAInlineRefTarget{sga3:XXIII:1.1:item:v}(v)] an element
    \(X_\alpha\in\Gamma(S,\mathfrak g^\alpha)^\times\) for every
    \(\alpha\in\Delta\), that is, an element
    \[
      u_\alpha=\exp_\alpha(X_\alpha)\in U_\alpha^\times(S)
      \qquad\text{for every }\alpha\in\Delta,
    \]
\end{enumerate}
satisfying condition (D1) of \SGARefLink{sga3:xxii:def:1.13}{Exposé~XXII, 1.13} (indeed, condition
(D2) of \textit{loc.\ cit.} is automatic).

Every split group admits a pinning; in particular, every reductive
group is locally pinnable for the étale topology.
